\section{Dataset Construction and Used Features}
\label{sec:dataset}

The dataset used in this project is based on the Capital Bikeshare system 
from the Washington D.C. area \cite{taweilo_kaggle,capital_system_data}. 
It contains historical trip records, station metadata, and daily weather 
data from May 2020 to August 2024. The raw trip table was aggregated into 
a station-day format: each row represents exactly one start station on one 
calendar day. This reduction aligns rental counts with daily weather 
information, produces a single well-defined target variable, and avoids 
mixing row granularities. After filtering to ensure identical station-day 
keys for both the lag and no-lag variants, the final dataset contains 
30,840 rows, split chronologically into 21,410 training, 4,705 validation, 
and 4,725 test rows (70/15/15).

\subsection{Target Variable and Distribution}

The task is a regression problem with a single numeric target: 
\texttt{total\_rentals}, the total number of bike trips started at a 
specific station on a given day. Intermediate count columns for 
member/casual users or bike types were removed from the predictor set 
because they are direct components of the same total. Keeping them would 
create a leakage-like shortcut where the model reconstructs demand instead 
of forecasting it.

\begin{table}[htbp]
\caption{Target variable distribution (30,840 station-day rows).}
\label{tab:target_dist}
\centering
\scriptsize
\begin{tabularx}{\columnwidth}{p{0.33\columnwidth} X}
\toprule
\textbf{Statistic} & \textbf{\texttt{total\_rentals}} \\
\midrule
Mean        & 73.77 rentals per station-day \\
Median      & 70 rentals \\
Std. dev.   & 39.75 rentals \\
Min / Max   & 1 / 377 rentals \\
\bottomrule
\end{tabularx}
\end{table}

The distribution shows a wide range of demand levels. This justifies 
using both RMSE and MAE as evaluation metrics, since MAE alone does not 
penalize large prediction errors sufficiently.

\subsection{Station Filtering and One-Hot Encoding}

We focused on the 20 most popular start stations by total trip count. 
Rare stations provide fewer repeated observations, which makes 
station-specific patterns unstable and weakens model evaluation. The 
trade-off is explicit: our models represent high-demand stations, not 
the full bikeshare system.

The station identity was originally stored as a numeric ID. Using this 
directly as a feature would mislead models into treating the ID as an 
ordered quantity. We therefore applied one-hot encoding, producing 20 
binary columns, one per station \cite{sklearn_ohe}.

\subsection{Feature Groups}

To make accurate predictions, the models need to understand seasonal 
patterns, weather conditions, and recent demand history. Before 
finalising the feature set, we analysed pairwise correlations 
(see Figure~\ref{fig:corr}) to remove near-duplicate weather variables 
such as \texttt{feelslike} and \texttt{temp}.

\begin{table}[htbp]
\caption{Input features used in the final modelling dataset.}
\label{tab:features}
\centering
\scriptsize
\begin{tabularx}{\columnwidth}{p{0.28\columnwidth} X}
\toprule
\textbf{Feature group} & \textbf{Variables and description} \\
\midrule
Station identity & 
20 one-hot binary columns from \texttt{start\_station\_id}. 
Captures demand differences between stations without implying 
a numeric order. \\
\addlinespace
Weather \& daylight & 
\texttt{tempmax}, \texttt{humidity}, \texttt{windspeed}, 
\texttt{precip}, \texttt{precipcover}, \texttt{cloudcover}, 
\texttt{snow}, \texttt{snowdepth}, \texttt{visibility}, 
\texttt{sealevelpressure}, \texttt{uvindex}, 
\texttt{sunset\_minutes} (daylight duration proxy). \\
\addlinespace
Calendar \& time & 
\texttt{month}, \texttt{weekday} (seasonality and commuting 
patterns), \texttt{year}, \texttt{time\_idx} (continuous 
chronological index for long-term trends). \\
\addlinespace
Temporal memory & 
\texttt{total\_rentals\_lag\_1} (yesterday's demand) and 
\texttt{total\_rentals\_lag\_7} (same weekday last week), 
both station-specific. \\
\bottomrule
\end{tabularx}
\end{table}

\begin{figure}[htbp]
\centering
\includegraphics[width=\linewidth]{predictor_correlation_heatmap.png}
\caption{Correlation matrix used during feature reduction. Strong 
blocks identified near-duplicate weather variables and supported 
removing redundant predictors before modelling.}
\label{fig:corr}
\end{figure}